%! Author = sriadityavedantam
%! Date = 3/25/26

\section{Axiom of Completness}

\lesson{2}{Friday 15 August 2025 9:10}{Day 2}

\begin{definition}[Ordering Property]
	Given $a,b\in\r$, either
	\begin{align*}
		a &> b\\
		a &= b\\
		a &< b.
	\end{align*}
\end{definition}

\begin{remark}
	Field $\r$ is complete, which means it has the least upper bound (l.u.b) property.
\end{remark}

\begin{definition}[l.u.b Property]
	If $S\subset\r$ and $S$ is bounded above, i.e.\ there is a $C\in\r$ such that $C\geq x$ for all $x\in S$, then there is some $\alpha\in\r$ which is the least upper bound for $S$.
	This means $\alpha\leq C$ for all upper bounds $C$ of $S$.
	\[
	\alpha \coloneqq \sup(S)
	\]
\end{definition}

\begin{corollary}
{The greatest lower bound property follows.}
	\[\beta \coloneqq \inf(S)\]
\end{corollary}

\begin{corollary}
{\[\inf(S) = -\sup\{-x : x\in S\}.\]}
	Prove this as an Exercise.
\end{corollary}

\begin{remark}
	This is a convincing proof for Dr. Fu, but for now we must still be able to prove it.
\end{remark}

\begin{corollary}
{The Archimedian Property of $\r$ states:
    \begin{displayquote}
        For any $x\in\r$, there exists $n\in\n$ such that $n > x$.
    \end{displayquote}}
	Suppose for the sake of contradiction, there is no $n > x$.
	Then $\n$ is bounded above.
	Let $\alpha \coloneqq\sup(\n)$.
	Since $\alpha - 1 < \alpha$, then $\alpha - 1$ cannot be an upper bound for $\n$, so there is a $m\in\n$ such that $\alpha - 1 < m$.
	Thus $\alpha < m + 1$, which contradicts that $\alpha$ is an upper bound.
\end{corollary}

\begin{remark}
	The direct statement of this corollary is:
	\begin{displayquote}
		If $x\in\r$, then there exists $n\in\n$ such that $n>x$.
	\end{displayquote}
	A contrapositive of this lemma is:
	\begin{displayquote}
		If there is no $n\in\n$ such that $n>x$, then $x\notin\r$.
	\end{displayquote}
\end{remark}

\begin{corollary}{
	For every real $\epsilon > 0$, there is an $n\in\n$ such that $\dfrac{1}{n} < \epsilon$.
}
	By the archimedian principle, there exists $n\in\n$ such that $n > \dfrac{1}{\epsilon}$ thus $\dfrac{1}{n} < \epsilon$.
\end{corollary}

\begin{remark}
	The direct statement of this corollary is:
	\begin{displayquote}
		If $\eps > 0$ is real, then there exists $n\in\n$ such that $\dfrac{1}{n} < \eps$.
	\end{displayquote}
	A contrapositive of this lemma is:
	\begin{displayquote}
		If $\dfrac{1}{n} \geq \eps$ for all $n\in\n$, then $\eps \leq 0$.
	\end{displayquote}
\end{remark}

\begin{lemma}[Ordering of Squares]{
	$a,b\in\r$, $a,b > 0$, $a < b$ if and only if $a^2 < b^2$.
}
	Prove this as an Exercise.
\end{lemma}

\begin{corollary}{
	There is a real number $x > 0$ such that $x^2 = 2$.
}
	Define $S \coloneqq \{y\in\r : y^2 < 2\}$.
	By the previous lemma, $S$ is bounded above.
	Suppose $\alpha\coloneqq\sup(S)$.
	Suppose $\alpha^2 < 2$.
	Then there is $\epsilon > 0$ such that $(\alpha + \epsilon)^2 < 2$.
	Then:
	\begin{align*}
		(\alpha + \epsilon)^2 &= \alpha^2 + 2\alpha\epsilon + \epsilon^2 < 2.
	\end{align*}
	Thus $\alpha + \epsilon \in S$, contradicting that $\alpha$ is an upper bound.
	Suppose $\alpha^2 > 2$.
	Then there is $\epsilon > 0$ such that $(\alpha - \epsilon)^2 > 2$.
	Then:
	\begin{align*}
		(\alpha - \epsilon)^2 &= \alpha^2 - 2\alpha\epsilon + \epsilon^2 > 2.
	\end{align*}
	Thus $\alpha - \epsilon$ is still an upper bound smaller than $\alpha$, contradicting minimality.
	Hence $\alpha^2 = 2$.
	Uniqueness follows from the ordering of squares.
\end{corollary}